\begin{textsample}{POS Dim 2 – 2010\_c5 – Score -1.00 – 2010\_c5/enade\_2010\_biomed\_6d2a65c171.txt}  \label{ex:f2_pos_326}
CAPÍTULO 5 DISTRIBUIÇÃO DOS CONCEITOS 5.1 PANORAMA NACIONAL DA DISTRIBUIÇÃO DOS CONCEITOS A Tabela 5.1 apresenta a quantidade e distribuição de cursos de Biomedicina participantes do ENADE/2010 por faixa de conceito e Grande Região.
Observando -se os dados da Tabela 5.1, nota -se que, dos 176 cursos participantes, 29,0 \% ( 51 cursos ) classificaram -se com conceito 2, o valor modal em âmbito nacional e para todas as regiões.
Os cursos sem conceito, por sua vez, representaram 34,1 \% do total nacional ( 60 cursos ).
A região Norte participou com dez cursos, 5,7 \% do total nacional.
Nesta região, o conceito 2 foi atribuído a quatro dos dez participantes, equivalente a 40,0 \% do total regional.
Um pouco menos de 1/3 dos cursos ( 30,0 \%, correspondendo a três cursos ) não receberam conceito.
Dois cursos receberam conceito 3, enquanto apenas um obteve conceito 4.
A região Nordeste participou com 27 cursos, correspondentes a 15,3 \% do total nacional.
Destes, oito cursos, 29,6 \% em termos regionais, obtiveram conceito 2, o conceito modal.
Tanto o conceito 1 quanto o conceito 3 tiveram, cada um, três cursos avaliados.
O conceito 4 foi atribuído a dois cursos ( 7,4 \% ), e o conceito 5, a apenas um curso ( 3,7 \% ).
Nessa região, 37,1 \% dos cursos ( dez ) ficaram sem conceito.
Na região Sudeste a predominância foi também do conceito 2 ( 28,4 \% ), correspondente a 25 dos 88 cursos participantes na região.
O conceito 1 foi atribuído a quatro cursos ( 4,5 \% ) e o conceito 3, a 22 ( 25,0 \% ).
O conceito 4 foi correspondeu a sete cursos ( 8,0 \% ) e o conceito 5, a um curso ( 1,1 \% ).
Do total de cursos da região Sudeste, 29 ( 33,0 \% ) ficaram sem conceito.
Dos 32 cursos participantes da região Sul, oito ( 25,0 \% ) obtiveram conceito 2, o conceito modal.
O conceito 1 foi atribuído a dois cursos ( 6,3 \% ) e o conceito 3, a sete ( 21,9 \% ).
Apenas um curso ( 3,1 \% ) recebeu conceito 4, enquanto três deles ( 9,4 \% ) receberam conceito 5.
Os 11 cursos remanescentes ( 34,3 \% ) ficaram sem conceito.
Na região Centro-Oeste, seis cursos receberam o conceito modal, já mencionado.
Três cursos ( 15,8 \% ) foram avaliados com conceito 3, dois ( 10,5 \% ) com conceito 4 e apenas um ( 5,3 \% ) com conceito 5.
Também foi observado que sete cursos ( 36,8 \% ) não receberam conceito. 5.2 CONCEITOS POR CATEGORIA ADMINISTRATIVA E POR GRANDE REGIÃO A Tabela 5.2 apresenta a distribuição dos cursos participantes do ENADE/2010 de Biomedicina, por Categoria Administrativa, de acordo com os conceitos por eles alcançados, segundo as Grandes Regiões brasileiras.
Dos 176 cursos participantes, 23 ( 13,1 \% ) eram ministrados em instituições públicas e 153 ( 86,9 \% ), em privadas.
De acordo com as informações da Tabela 5.2, as instituições públicas foram as que apresentaram, em termos nacionais, a maioria dos cursos com conceito 5 ( cinco dos seis cursos foram classificados com esse conceito ).
Dos 23 cursos participantes de IES públicas, o conceito 4 foi atribuído a seis cursos, o conceito modal.
Classificaram -se, ainda, dois cursos com conceito 1, um com conceito 2 e três com conceito 3.
Seis cursos ficaram sem conceito nessa Categoria Administrativa.
Na rede privada, o conceito modal foi 2 com 50 cursos dos 153 desse tipo de IES.
O conceito 3 foi atribuído a 34 cursos.
Dois conceitos – 1 e 4 – tiveram sete cursos avaliados em cada um.
Apenas um curso obteve conceito 5.
Nesta Categoria, 54 cursos ficaram sem conceito.
Na análise por região, observa -se que, na região Norte, as instituições privadas participaram com nove cursos, dos quais três ficaram sem conceito.
O conceito modal para as instituições privadas na região Norte foi 2 ( com quatro cursos ).
Os demais foram avaliados com os conceitos 3 e 4 ( um curso cada ).
Três cursos dentre os nove não receberam conceito.
As instituições públicas participaram com 1 curso apenas, e este recebeu o conceito modal, 3.
Na região Nordeste, a rede privada concentrou 23 dos 27 cursos participantes, 85,2 \% do total da região.
Entre os que receberam conceito nesta Categoria Administrativa, predominaram os cursos com conceito 5 ( nove cursos ).
Em sequência, oito cursos obtiveram conceito 1.
Os demais foram avaliados com conceito 2 ( três cursos ), ou ficaram sem conceito ( três cursos ).
As instituições públicas nesta região participaram com quatro cursos, dos quais dois foram avaliados com conceito 3, um com conceito 4 e um outro com conceito 5.
Na região Sudeste, a rede privada contou com 88,6 \% dos cursos participantes, \textbf{totalizando} 78.
Nesta categoria na região, o conceito modal foi 2 ( 24 cursos ).
Os demais foram avaliados com conceito 1 ( três cursos ), conceito 3 ( 21 cursos ) e conceito 4 ( quatro cursos ).
Na combinação de Categoria Administrativa e Grande Região, 26 cursos não receberam conceito.
Entre os dez cursos ( 11,4 \% ) em instituições públicas na região, a categoria modal foi a que recebeu conceito 4 ( três cursos ).
Os demais tiveram apenas um curso cada distribuídos pelos conceitos 1, 2, 3 e 5.
Nesta categoria na região Sudeste, três cursos ficaram sem conceito.
As instituições privadas concentraram 27 dos 32 cursos participantes da região Sul, 84,4 \% do total regional.
Desses, oito ficaram classificados no conceito modal, 2.
Os demais foram avaliados com conceito 1 ( um curso ), 3 ( sete cursos ) e 4 ( apenas um deles ).
Na combinação de Categoria Administrativa e Grande Região, dez cursos não receberam conceito.
As instituições públicas nesta região participaram com cinco cursos ( 15,6 \% ), dos quais três receberam conceito modal, 5.
Dos demais, um curso foi avaliado com conceito 1 e um curso ficou sem conceito.
Na região Centro-Oeste, 16 dos 19 cursos participantes eram de instituições privadas ( 84,2 \% em termos regionais ).
Destes, seis integram o conceito modal, 2.
Outros seis cursos não receberam conceito.
Simultaneamente, dois cursos receberam o conceito 3, um recebeu conceito 4 e outro conceito 5.
Na rede pública, dos três cursos ( 15,8 \% ) na região Centro-Oeste, um ficou sem conceito e os outros dois ficaram com os conceitos 3 e 4. 5.3 CONCEITOS POR ORGANIZAÇÃO ACADÊMICA E POR GRANDE REGIÃO Na Tabela 5.3 encontra -se a distribuição dos conceitos atribuídos aos cursos participantes do ENADE/2010 na Área de Biomedicina, por Organização Acadêmica, segundo as Grandes Regiões brasileiras.
Dos 176 cursos participantes, 70 eram oferecidos em Universidades, 30 em Centros Universitários e os demais 76 em Faculdades.
De acordo com os dados apresentados, todos os cursos avaliados com conceito 5 eram vinculados a Universidades ( seis dos seis cursos com este conceito ).
Esse tipo de Organização Acadêmica teve conceito 3 como modal, com 19 cursos.
Os demais foram avaliados com conceito 1 ( cinco cursos ), conceito 2 ( 12 cursos ) e conceito 4 ( 11 cursos ), além dos 17 cursos que ficaram sem conceito.
Entre os cursos em Centros Universitários, o conceito modal ficou entre 2 e 3, com dez cursos cada, e outros sete ficaram sem conceito.
Os demais cursos, neste tipo de Organização Acadêmica, receberam conceito 1 ( um curso ) e conceito 4 ( dois deles ).
Nas Faculdades, dos 76 cursos, 36 ficaram sem conceito e 29 receberam o conceito modal, 2, além daqueles que receberam conceito 1 ( três cursos ) e conceito 3 ( oito cursos ).
Considerando -se separadamente as regiões brasileiras, verifica -se que, na região Norte, as Universidades concentraram dois dos dez cursos, sendo que um deles foi avaliado com conceito 3, e o outro, com conceito 4.
Os Centros Universitários também foram representadas por dois cursos, sendo um sem conceito e o outro com conceito 2.
As Faculdades participaram com seis cursos na região Norte.
Destes, dois ficaram sem conceito e os quatro demais receberam conceito 2 ( três cursos – o conceito modal ) e conceito 3 ( um curso ).
Na região Nordeste, as Universidades participaram com cinco dos 27 cursos na Área de Biomedicina da região.
Entre os cursos em Universidades no Nordeste, um ficou sem conceito.
O conceito modal foi 4 ( com dois cursos ).
Os demais dois cursos receberam conceito 2 ( um curso ) e conceito 5 ( também um ).
Os Centros Universitários na região Nordeste contaram com um único curso que recebeu conceito modal, 2.
As Faculdades foram representadas por 21 cursos, dos quais nove não receberam conceito.
O conceito modal foi 2, com seis cursos assim avaliados.
Além disso, três cursos receberam conceito 1 e outros três conceito 3.
Na região Sudeste, as Universidades concentraram 41 dos 88 cursos da região.
Entre os cursos nesse tipo de Organização Acadêmica no Sudeste, o conceito modal foi 3, com 15 cursos e 11 ficaram sem conceito.
Os demais cursos receberam conceito 1 ( três cursos ), conceito 2 ( seis cursos ), conceito 4 ( cinco cursos ) e conceito 5 ( apenas um curso ).
Os Centros Universitários participaram com 19 cursos na região Sudeste, dos quais seis obtiveram o conceito modal, 3, e cinco ficaram sem conceito.
Os demais receberam conceito 1 ( um curso ), conceito 2 ( cinco cursos ) e conceito 4 ( dois cursos ).
As Faculdades foram representadas por 28 cursos na região Sudeste, que se distribuíram nos conceitos 2 ( 14 cursos, o conceito modal ) e 3 ( um curso ).
Os 13 cursos restantes não receberam conceito.
Dos 32 cursos da região Sul, 16 eram de Universidades, para os quais o conceito modal foi 2 ( com quatro cursos ).
Nesse tipo de organização, quatro cursos ficaram sem conceito e os demais receberam conceito 1 ( dois cursos ), conceito 3 ( dois cursos ), 4 ( apenas um curso ) e conceito 5 ( três deles ).
Os Centros Universitários tiveram cinco cursos participantes, dos quais um ficou sem conceito e os outros quatro ficaram igualmente divididos entre os conceitos 2 e 3.
Foram 11 os cursos vinculados a Faculdades na região Sul, e destes, seis ficaram sem conceito.
Os cinco restantes receberam conceito 2 ( dois cursos ) e conceito 3, o conceito modal ( três cursos ).
Na região Centro-Oeste, seis dos 19 cursos eram de Universidades.
Nesse tipo de organização, um curso ficou sem conceito e o conceito modal foi 4, com dois cursos.
Os três cursos remanescentes foram avaliados, cada um, nos conceitos 2, 3 e 5.
Os Centros Universitários na região Centro-Oeste contaram com três cursos, sendo que um recebeu conceito 2 e os outros dois receberam conceito modal, 3.
Dos dez cursos em Faculdades da região Centro-Oeste, seis ficaram sem conceito e os quatro restantes foram avaliados com conceito 2.

% matched lemmas: totalizar
\end{textsample}
